% =====================================================================
%  Preamble for "Amber and 6G Ambient IoT: A Teaching Text"
% =====================================================================
\usepackage[T1]{fontenc}
\usepackage[utf8]{inputenc}
\usepackage{mathpazo}                 % Palatino text and math
\usepackage[scaled=0.92]{helvet}      % sans-serif for headings / labels
\renewcommand{\ttdefault}{lmtt}       % Latin Modern Typewriter for code
\usepackage{microtype}
\usepackage[a4paper,margin=2.5cm,headheight=14pt]{geometry}
\usepackage{setspace}
\setstretch{1.07}
\emergencystretch=2em

% ---- mathematics -------------------------------------------------------
\usepackage{amsmath,amssymb,mathtools,bm}
\usepackage{siunitx}
\sisetup{detect-all, per-mode=symbol, range-phrase={--},
         group-minimum-digits=4, input-ignore={,}, input-decimal-markers={.}}
\DeclareSIUnit{\dBm}{dBm}
\DeclareSIUnit{\dB}{dB}
\DeclareSIUnit{\dBi}{dBi}
\DeclareSIUnit{\mAh}{mAh}
\DeclareSIUnit{\ms}{ms}
\DeclareMathOperator{\sinc}{sinc}
\DeclareMathOperator*{\argmax}{arg\,max}
\newcommand{\E}{\mathbb{E}}
\newcommand{\e}{\mathrm{e}}
\newcommand{\dd}{\,\mathrm{d}}
\newcommand{\Rs}{R_{\mathrm{s}}}          % series resistance
\newcommand{\Rl}{R_{\mathrm{l}}}          % leakage resistance
\newcommand{\Vhigh}{V_{\mathrm{high}}}
\newcommand{\Vlow}{V_{\mathrm{low}}}
\newcommand{\Vmax}{V_{\mathrm{max}}}
\newcommand{\Vs}{V_{\mathrm{src}}}
\newcommand{\Pdbm}{P_{\mathrm{dBm}}}

% ---- tables, figures ---------------------------------------------------
\usepackage{booktabs,tabularx,multirow,array,makecell,longtable}
\usepackage{graphicx}
\graphicspath{{figures/}}
\usepackage[dvipsnames,table]{xcolor}
\usepackage{caption,subcaption}
\captionsetup{font=small,labelfont={bf,sf},format=hang,margin=0.5em}
\usepackage{float}
\usepackage{enumitem}
\setlist{itemsep=2pt,topsep=4pt,parsep=0pt}

% ---- palette (matches the Python figures) ------------------------------
\definecolor{ablue}{HTML}{1F5FA8}
\definecolor{aorange}{HTML}{E07B00}
\definecolor{agreen}{HTML}{2A8F4A}
\definecolor{ared}{HTML}{C9302C}
\definecolor{apurple}{HTML}{7B4FA6}
\definecolor{ateal}{HTML}{16808A}
\definecolor{agrey}{HTML}{6B6B6B}
\definecolor{agold}{HTML}{F2C14E}
\definecolor{lightblue}{HTML}{EAF1FA}
\definecolor{lightorange}{HTML}{FFF3E3}
\definecolor{lightgreen}{HTML}{EAF6EE}
\definecolor{lightred}{HTML}{FBEAEA}
\definecolor{lightgrey}{HTML}{F3F3F3}

% ---- TikZ --------------------------------------------------------------
\usepackage{tikz}
\usepackage[siunitx]{circuitikz}
\usetikzlibrary{arrows.meta,positioning,shapes.geometric,fit,calc,
  decorations.pathreplacing,decorations.markings,automata,backgrounds,
  patterns,matrix,chains}
\tikzset{
  >=Stealth,
  box/.style={draw,rounded corners=2pt,align=center,inner sep=4pt,
              font=\small\sffamily,fill=white},
  mod/.style={box,fill=lightblue,draw=ablue,minimum height=0.9cm},
  ext/.style={box,fill=lightgrey,draw=agrey},
  st/.style={draw,rounded corners=3pt,align=center,inner sep=3pt,
             font=\footnotesize\sffamily,minimum height=0.75cm},
  lab/.style={font=\scriptsize\sffamily,inner sep=1.5pt},
  every node/.style={font=\small}
}

% ---- algorithms / listings ---------------------------------------------
\usepackage{algorithm}
\usepackage[noend]{algpseudocode}
\algrenewcommand\algorithmicrequire{\textbf{Input:}}
\algrenewcommand\algorithmicensure{\textbf{Output:}}
\usepackage{listings}
\lstset{basicstyle=\ttfamily\footnotesize,columns=fullflexible,
  keepspaces=true,breaklines=true,frame=single,rulecolor=\color{agrey!60},
  backgroundcolor=\color{lightgrey},language=Python,showstringspaces=false,
  keywordstyle=\color{ablue}\bfseries,commentstyle=\color{agrey}\itshape,
  stringstyle=\color{agreen},xleftmargin=0.5em,xrightmargin=0.5em,
  aboveskip=6pt,belowskip=6pt}

% ---- teaching callouts -------------------------------------------------
\usepackage[most]{tcolorbox}
\tcbset{boxsep=3pt,left=6pt,right=6pt,top=4pt,bottom=4pt,
        fonttitle=\bfseries\sffamily\small,arc=2pt,boxrule=0.6pt,
        breakable,enhanced jigsaw}
\newtcolorbox{keyidea}[1][Key idea]{colback=lightblue,colframe=ablue,
  title={#1},attach boxed title to top left={xshift=8pt,yshift=-2mm},
  boxed title style={colback=ablue,arc=2pt}}
\newtcolorbox{defbox}[1][Definition]{colback=white,colframe=agrey,
  title={#1},fonttitle=\bfseries\sffamily\small,colbacktitle=lightgrey,
  coltitle=black}
\newtcolorbox{example}[1][Worked example]{colback=lightgreen,colframe=agreen,
  title={#1},attach boxed title to top left={xshift=8pt,yshift=-2mm},
  boxed title style={colback=agreen,arc=2pt}}
\newtcolorbox{reality}[1][Code reality check]{colback=lightorange,colframe=aorange,
  title={#1},attach boxed title to top left={xshift=8pt,yshift=-2mm},
  boxed title style={colback=aorange,arc=2pt}}
\newtcolorbox{intuition}[1][Intuition]{colback=white,colframe=apurple,
  title={#1},attach boxed title to top left={xshift=8pt,yshift=-2mm},
  boxed title style={colback=apurple,arc=2pt}}
\newtcolorbox{codemap}[1][Where this lives in the code]{colback=lightgrey,
  colframe=agrey,title={#1},fonttitle=\bfseries\sffamily\small}
\newtcolorbox{checkpoint}[1][Checkpoint]{colback=lightred,colframe=ared,
  title={#1},attach boxed title to top left={xshift=8pt,yshift=-2mm},
  boxed title style={colback=ared,arc=2pt}}

% numbered exercises (solutions are collected in an appendix)
\newcounter{exercise}[chapter]
\renewcommand{\theexercise}{\thechapter.\arabic{exercise}}
\newtcolorbox[use counter=exercise]{exercise}[1][]{colback=white,colframe=ablue!70,
  title={Exercise~\theexercise}, fonttitle=\bfseries\sffamily\small,
  colbacktitle=ablue!12,coltitle=black,#1}

% ---- code identifiers --------------------------------------------------
\newcommand{\code}[1]{\texttt{\detokenize{#1}}}
\newcommand{\file}[1]{\texttt{\detokenize{#1}}}
\usepackage{seqsplit}
\newcommand{\codew}[1]{\texttt{\expandafter\seqsplit\expandafter{\detokenize{#1}}}}  % breakable, for table cells

% ---- headings ----------------------------------------------------------
\usepackage{titlesec}
\titleformat{\chapter}[display]
  {\sffamily\bfseries\Large\color{ablue}}
  {\filleft\fontsize{54}{54}\selectfont\thechapter}
  {4pt}{\filleft\Huge\color{black}}[\vspace{2pt}\titlerule]
\titleformat{name=\chapter,numberless}[display]
  {\sffamily\bfseries\Large\color{ablue}}{}{0pt}{\filleft\Huge\color{black}}[\titlerule]
\titlespacing*{\chapter}{0pt}{-10pt}{26pt}
\titleformat{\section}{\sffamily\bfseries\Large\color{ablue}}{\thesection}{0.7em}{}
\titleformat{\subsection}{\sffamily\bfseries\large}{\thesubsection}{0.7em}{}
\titleformat{\subsubsection}{\sffamily\bfseries\normalsize}{\thesubsubsection}{0.7em}{}
\titleformat{\paragraph}[runin]{\sffamily\bfseries\small}{}{0pt}{}
\titleformat{\part}[display]
  {\sffamily\bfseries\Huge\color{ablue}}{\filleft\partname~\thepart}{10pt}{\filleft}

\usepackage{fancyhdr}
\pagestyle{fancy}
\fancyhf{}
\fancyhead[LE]{\sffamily\small\nouppercase{\leftmark}}
\fancyhead[RO]{\sffamily\small\nouppercase{\rightmark}}
\fancyfoot[C]{\sffamily\small\thepage}
\renewcommand{\headrulewidth}{0.3pt}
\fancypagestyle{plain}{\fancyhf{}\fancyfoot[C]{\sffamily\small\thepage}%
  \renewcommand{\headrulewidth}{0pt}}

% ---- bibliography, links -----------------------------------------------
\usepackage[numbers,sort&compress]{natbib}
\usepackage[hidelinks,pdfusetitle,bookmarksnumbered]{hyperref}
\hypersetup{colorlinks=true,linkcolor=ablue,citecolor=agreen,urlcolor=ablue}
\usepackage[nameinlink,capitalise]{cleveref}
\usepackage{appendix}
\setcounter{tocdepth}{2}
\setcounter{secnumdepth}{3}
\usepackage{tocloft}
\renewcommand{\cftchapfont}{\sffamily\bfseries}
\renewcommand{\cftsecfont}{\sffamily}
\renewcommand{\cftchappagefont}{\sffamily\bfseries}
\renewcommand{\cftsecpagefont}{\sffamily}
\hfuzz=2pt
